\chapter{내부 표현을 읽는 방법과 그 한계}
\label{chap:phase4}

\section{결론부터: 읽을 수 있다는 사실은 사용된다는 뜻이 아니다}

활성값에서 어떤 변수를 예측할 수 있다는 결과는 그 변수에 관한 정보가 활성값에 존재한다는 증거다.
그러나 원 모델이 같은 판별 경계나 같은 방향을 실제 출력 계산에 사용한다는 결론은 별도의 인과적 증거가 없으면 성립하지 않는다.
이 장의 핵심 구분은 다음 세 문장이다.

\begin{enumerate}
  \item 변수 $y$를 활성값 $\vect{h}$에서 \emph{복원할 수 있다}.
  \item $y$에 대응한다고 해석한 방향을 바꾸면 모델 행동도 \emph{바뀐다}.
  \item 원 모델이 자연스러운 입력에서 그 방향을 이용하는 \emph{메커니즘}을 안다.
\end{enumerate}

첫 문장은 관찰적이고, 둘째 문장은 개입적이며, 셋째 문장은 구성요소와 경로를 포함한 설명이다.
앞 문장이 참이라고 해서 다음 문장이 자동으로 참이 되지는 않는다.

\section{분산 표현과 특징 방향}

분산 표현은 의미 변수가 좌표 하나가 아니라 여러 좌표의 결합과 관계를 맺을 수 있다는 관점이다.
layer $\ell$, position $t$의 residual state를 $\vect{h}^{(\ell)}_t\in\R^d$라고 할 때, ``개체'', ``시제'', ``진실성'' 같은 변수가 좌표 한 개에만 저장되어야 할 이유는 없다.
한 좌표도 여러 입력 조건에서 서로 다른 변수와 상관될 수 있다.

선형 readout은
\begin{equation}
  s(\vect{h})=\vect{w}^{\mathsf T}\vect{h}+b
  \label{eq:linear-readout}
\end{equation}
로 score를 만든다.
binary classification에서는 $s>0$인지로 class를 나눌 수 있다.
$\vect{w}$를 feature direction이라고 부를 때는 다음 차이를 유지해야 한다.

\begin{itemize}
  \item 좌표 또는 neuron은 선택한 basis의 축 하나다.
  \item 방향은 여러 좌표의 선형결합일 수 있다.
  \item subspace는 한 개 이상의 독립 방향이 만드는 부분공간이다.
  \item feature는 연구자가 정한 의미 변수 또는 학습된 dictionary element를 가리킬 수 있으므로, 논문마다 조작적 정의를 확인해야 한다.
\end{itemize}

개별 neuron의 의미 해석은 모델이 사용하는 basis에 의존한다.
invertible matrix $\mat{A}$로 좌표를 $\vect{h}'=\mat{A}\vect{h}$로 바꾸고 다음 layer의 weight를 그에 맞게 보정하면 같은 함수를 표현할 수 있지만 개별 좌표의 값은 달라진다.
따라서 ``neuron 1723이 개념 그 자체다''라는 주장은 basis와 입력 분포를 명시하지 않으면 강한 주장이다.
반면 특정 모델의 고정된 basis에서 그 neuron의 활성 조건을 기술하는 것은 검증 가능한 경험적 주장이다.

\begin{interpretbox}{분산 표현의 최소 의미}
분산 표현은 모든 정보가 모든 좌표에 균등하게 흩어져 있다는 뜻이 아니다.
어떤 변수는 소수의 방향이나 neuron에서 더 쉽게 복원될 수 있다.
이 용어는 의미 변수와 parameter 좌표 사이에 필연적인 일대일 대응이 없다는 점을 가리킨다.
\end{interpretbox}

\section{표현 공간의 기하}

표현 기하 연구는 활성값 사이의 거리, 각도, 군집, 저차원 구조와 변환 관계를 분석한다.
두 벡터의 cosine similarity는
\begin{equation}
  \cos(\vect{x},\vect{y})=
  \frac{\vect{x}^{\mathsf T}\vect{y}}{\lVert\vect{x}\rVert_2
  \lVert\vect{y}\rVert_2}
\end{equation}
다.
이는 방향의 일치도를 재지만 norm 차이를 버린다.
residual norm이 logit이나 다음 normalization에 영향을 주는 모델에서는 cosine만으로 계산 효과를 판단할 수 없다.

주성분분석은 표본 분산이 큰 방향을 찾지만 그 방향을 인간 개념으로 보장하지 않는다.
데이터 행렬 $\mat{H}\in\R^{n\times d}$를 평균 중심화하고 SVD를 적용하면 이 방향을 계산할 수 있다.
label에 따라 평균 차이
\begin{equation}
  \vect{v}=\E[\vect{h}\mid y=1]-\E[\vect{h}\mid y=0]
\end{equation}
를 계산하는 방법도 있다.
이 방향은 표본 분포, layer, position과 prompt 형식에 의존한다.

일부 연구는 개념이나 관계가 선형 구조로 나타난다는 ``선형 표현 가설''을 이론과 실험으로 분석한다 \citep{park2024linear}.
이 문헌은 특정 정의와 모델에서 선형 구조를 지지하는 결과를 제공하지만, 모든 의미 변수가 한 개의 전역 선형 방향으로 표현된다는 정리를 제공하지 않는다.

\begin{limitbox}
군집 그림은 차원 축소 방법과 hyperparameter에 민감할 수 있다.
t-SNE나 UMAP의 2차원 배치에서 떨어져 보이는 두 군집이 원래 공간에서도 같은 방식으로 분리된다고 단정해서는 안 된다.
원 차원의 거리, held-out 분류와 개입 결과를 함께 확인해야 한다.
\end{limitbox}

\section{선형 probe와 control task}

probe는 고정된 모델의 활성값을 입력으로 받아 label을 예측하도록 별도 모델을 학습하는 방법이다.
선형 probe의 대표 목적함수는
\begin{equation}
  \hat{\vect{w}},\hat b
  =\argmin_{\vect{w},b}
  \sum_{i=1}^{n}\operatorname{CE}
  \bigl(y_i,\sigma(\vect{w}^{\mathsf T}\vect{h}_i+b)\bigr)
  +\lambda\lVert\vect{w}\rVert_2^2
\end{equation}
다.
평가에는 probe를 학습하지 않은 held-out sample을 사용해야 한다.
\index{probe@선형 probe}

높은 정확도는 다음 대안들 가운데 무엇이 맞는지 혼자 결정하지 못한다.

\begin{itemize}
  \item 원 모델이 행동을 만들 때 실제로 사용하는 변수가 복원되었다.
  \item 입력의 표면적 confound가 label과 함께 복원되었다.
  \item 고차원 표현에서 probe가 우연한 label mapping을 암기했다.
  \item 해당 정보는 존재하지만 downstream 계산이 무시한다.
\end{itemize}

Hewitt와 Liang은 언어 구조 probe가 representation의 구조와 probe 자체의 암기 능력을 혼동할 수 있음을 지적하고, 무작위 control label에 대한 성능과 실제 과제 성능의 차이를 selectivity로 평가했다 \citep{hewitt2019probes}.
실제 연구에서는 train/test 분리뿐 아니라 lexical overlap, prompt template, token position, class balance와 probe capacity를 통제해야 한다.

\begin{studybox}[probe 결과를 기록하는 양식]
재현 가능한 probe 결과에는 모델부터 결론의 증거 수준까지 전체 분석 명세가 필요하다.
모델과 checkpoint, layer와 hook 위치, 사용한 token position, label의 정의, train/validation/test 분리 단위, probe class와 regularization, baseline, control task, metric과 confidence interval을 기록한다.
마지막으로 ``정보 존재''와 ``인과적 사용'' 중 어디까지 결론낼 수 있는지 한 문장으로 제한한다.
\end{studybox}

\section{logit lens와 tuned lens}

logit lens는 중간 residual state를 최종 vocabulary 공간으로 투영하여 각 layer에서 어떤 token이 선호되는지 관찰한다.
가장 단순한 형태는
\begin{equation}
  \vect{z}^{(\ell)}_t=
  \mat{W}_U\operatorname{Norm}_f(\vect{h}^{(\ell)}_t)+\vect{b}_U
  \label{eq:logit-lens}
\end{equation}
다.
최종 layer에 맞춰 학습된 normalization과 unembedding을 중간 layer에 적용하는 진단 도구다.
원 모델이 실제로 중간마다 LM head를 실행한다는 뜻은 아니다.
\index{logit lens@logit lens}

중간 residual distribution은 최종 residual distribution과 다를 수 있다.
tuned lens는 layer별 affine translator $\mat{A}_{\ell},\vect{c}_{\ell}$를 학습하여
\begin{equation}
  \vect{z}^{(\ell)}_t=
  \mat{W}_U\operatorname{Norm}_f
  (\mat{A}_{\ell}\vect{h}^{(\ell)}_t+\vect{c}_{\ell})+\vect{b}_U
\end{equation}
로 예측을 읽는다 \citep{belrose2023tunedlens}.
이는 중간 예측을 더 잘 맞추기 위한 별도 모델이며, translator의 계산이 원 Transformer 안에 존재한다는 증거는 아니다.

logit lens에서 정답 token의 rank가 layer에 따라 오르는 현상은 유용한 관찰이다.
그러나 그 token이 실제 생성에 필요했던 경로를 확인하려면 activation patching이나 path-level 개입 같은 추가 실험이 필요하다.

\section{중첩과 다의적 neuron}

중첩(superposition)은 서로 직교하지 않는 여러 feature 방향을 같은 표현 공간에 배치하는 가설이다.
차원 $d$보다 많은 잠재 feature를 표현해야 하고 feature가 희소하게 나타날 때 이런 배치가 가능하다.
toy model 연구는 sparsity와 feature importance가 달라질 때 비직교 feature 표현과 다의적 neuron이 나타나는 조건을 구성적으로 보였다 \citep{elhage2022superposition}.
\index{superposition@중첩(superposition)}

toy model에서 중첩을 구성한 결과와 자연 LLM에서 중첩을 식별한 결과는 서로 다른 명제다.

\begin{enumerate}
  \item 설계된 toy model에서 최적화가 중첩 해를 만든다.
  \item 자연어로 학습한 대형 LLM의 특정 activation이 같은 이유로 다의적이다.
\end{enumerate}

첫째는 통제된 모델에서 직접 조사할 수 있다.
둘째는 자연 모델의 feature 정의와 측정법을 추가로 요구한다.
한 neuron이 서로 무관해 보이는 여러 문맥에서 활성화된다는 관찰은 다의성의 증거가 될 수 있지만, 그 원인이 오직 중첩이라고 유일하게 식별하지는 못한다.

\section{희소 오토인코더}

희소 오토인코더(SAE)는 activation $\vect{x}\in\R^d$를 더 큰 feature 공간의 희소한 code $\vect{f}\in\R^m$으로 바꾸고 다시 복원한다.
단순한 한 형태는
\begin{align}
  \vect{f}&=\operatorname{ReLU}
  \bigl(\mat{W}_{\mathrm{enc}}(\vect{x}-\vect{b}_{\mathrm{dec}})
  +\vect{b}_{\mathrm{enc}}\bigr),\\
  \hat{\vect{x}}&=\mat{W}_{\mathrm{dec}}\vect{f}+\vect{b}_{\mathrm{dec}},\\
  \mathcal{L}_{\mathrm{SAE}}&=\lVert\vect{x}-\hat{\vect{x}}\rVert_2^2
  +\lambda\lVert\vect{f}\rVert_1
  \label{eq:sae}
\end{align}
로 쓸 수 있다.
$m>d$인 overcomplete dictionary를 사용할 수 있다.
decoder의 column은 activation 공간에 쓰는 방향이고, code coordinate는 해당 방향의 활성 세기다.
\index{sparse autoencoder@희소 오토인코더(SAE)}

희소 dictionary가 사람이 해석하기 쉬운 feature를 줄 수 있다는 대규모 사례 연구가 보고되었다 \citep{bricken2023monosemantic,templeton2024scaling}.
그러나 reconstruction loss와 sparsity penalty는 의미론적 유일성을 직접 최적화하지 않는다.
dictionary는 random seed, 데이터, sparsity 강도와 architecture에 따라 달라질 수 있으며, reconstruction error가 원 모델 계산의 일부를 잃는다.

SAEBench는 여러 SAE architecture를 여덟 종류의 metric으로 비교했고, 일반적인 proxy metric의 향상이 feature disentanglement나 practical downstream 성능의 향상으로 안정적으로 이어지지 않는 결과를 보고했다 \citep{karvonen2025saebench}.
따라서 낮은 reconstruction error 또는 높은 sparsity 하나만으로 ``해석가능성이 좋아졌다''고 결론내리면 안 된다.

\begin{limitbox}
SAE feature에 자연어 이름을 붙이는 일은 가설 생성이다.
feature가 활성되는 positive와 negative example, 유사 feature, activation distribution, decoder 방향 개입, downstream 효과와 reconstruction error 경로를 함께 조사해야 한다.
자동 설명 모델의 문장이 그 feature의 완전한 정의가 되는 것은 아니다.
\end{limitbox}

\section{표현 분석의 최소 증거 기준}

\begin{table}[htbp]
\centering
\caption{표현에 관한 주장과 필요한 추가 증거}
\label{tab:representation-evidence}
\begin{tabularx}{\textwidth}{@{}p{0.19\textwidth}YY@{}}
\toprule
주장 & 최소 관찰 & 아직 남는 질문 \\
\midrule
상관이 있다 & 독립 표본에서 activation과 label의 관계가 재현된다. & confound인가? \\
선형 복원이 된다 & held-out probe가 적절한 baseline과 control을 넘는다. & 원 모델이 사용하는가? \\
방향이 행동을 제어한다 & 크기와 부작용을 기록한 intervention이 행동을 바꾼다. & 자연 입력에서 같은 경로를 쓰는가? \\
구성요소가 필요하다 & 적절한 ablation에서 target metric이 나빠진다. & 중복 경로나 self-repair가 있는가? \\
메커니즘을 설명한다 & 경로 가설이 여러 반사실 입력과 개입 결과를 예측한다. & 분포와 모델을 넘어 일반화하는가? \\
\bottomrule
\end{tabularx}
\end{table}

표현 분석의 산출물은 결론문보다 재현 가능한 object여야 한다.
activation을 뽑은 정확한 hook, tensor shape, 표본 ID, 전처리, probe weight, metric과 seed를 보존해야 다음 장의 회로 분석과 인과 개입으로 이어갈 수 있다.

\section{Phase 4 학습 점검}

\begin{studybox}
\begin{enumerate}
  \item 같은 의미 label에 대해 neuron 하나, mean-difference direction과 regularized linear probe를 비교하고 held-out 성능을 기록한다.
  \item label을 무작위로 섞은 control task와 lexical split을 추가하여 selectivity를 평가한다.
  \item 각 layer의 logit lens와 tuned lens 분포를 비교하고, 어느 계산이 원 모델 밖에서 학습되었는지 표시한다.
  \item SAE 하나를 학습하여 reconstruction error, 평균 활성 feature 수, dead feature 비율과 intervention 결과를 함께 기록한다.
  \item 모든 그림과 표의 결론을 ``복원'', ``개입'', ``메커니즘'' 가운데 하나로 분류한다.
\end{enumerate}
\end{studybox}

\begin{factbox}
Phase 4의 종료 조건은 activation에 이름을 많이 붙이는 것이 아니다.
무엇이 관찰되었고, 어떤 대안 설명이 남으며, 다음 인과 실험이 무엇인지를 정확히 말할 수 있어야 한다.
\end{factbox}
